% Generated by scripts/blueprint_gen.py from docs/formalization-summary.tex.
\chapter{Introduction}\label{chap:intro}
\section{The theorem}

\paragraph{Sources.} Block \cite[Thm~4.1.3]{B}; Antweiler \cite[Thm~7.2.20 and Thm~7.4.3]{A};
\cite{CHL} for the comparison theorem in the abstract setting. For the analytic lemmas:
\cite{Hor, GH, DK, BSW}; for twisted sheaves: \cite{Cal, Toda}; for the triangulated structure:
\cite{BK}. The formalization is built on Mathlib \cite{Mathlib}.

\paragraph{Informal statement.} Let $E$ be a finite-dimensional complex normed space and $X$ a
compact Hausdorff complex manifold modelled on $E$. Let $A={\mathcal{A}}_X(X)={\mathcal{A}}^{0,\bullet}(X)$ be the
Dolbeault algebra of $X$ (smooth $(0,q)$-forms with ${\bar\partial}$ and wedge product) and ${\mathcal{O}}_X$ the
sheaf of holomorphic functions. Then the functor
\[ J\colon {\mathrm{Ho}}(\mathcal P_A)\longrightarrow D({\mathcal{O}}_X),\qquad M\longmapsto \alpha(M), \]
from the homotopy category of cohesive modules over $A$ to the derived category of sheaves of
${\mathcal{O}}_X$-modules, is full and faithful, and its essential image is ${D_{\mathrm{perf}}}(X,{\mathcal{O}}_X)$. In other
words, $J$ corestricts to an equivalence ${\mathrm{Ho}}(\mathcal P_A)\simeq{D_{\mathrm{perf}}}(X)$.

\paragraph{History of the statement.} The project first proved the theorem under two analytic
hypotheses on $X$:
\begin{itemize}
\item[(H1)] \emph{Dolbeault--Grothendieck lemma}: ${\bar\partial}$-closed $(0,q{+}1)$-forms are locally
${\bar\partial}$-exact;
\item[(H2)] \emph{local gauge lemma, weak form}: every cohesive module over $A$ is, near each
point, homotopy equivalent to one coming from a bounded complex of finitely generated projective
${\mathcal{O}}$-modules.
\end{itemize}
Both are true on every complex manifold (H1 is classical; H2 is the content of the proof of
\cite[Lemma~4.1.5]{B}), and both are now theorems of the project (\S\ref{chap:Z4}, \S\ref{chap:Z5}). The
conditional statement is kept, as a corollary, under the name \texttt{MainStatementConditional}.

\subsection{The Lean statement}
The statement lives in \texttt{Cohesivemodules/Statements.lean}, which imports only files marked
\texttt{-- FROZEN} (definitions checked in excerpts at a review gate, which may not be changed)
and Mathlib:
\begin{verbatim}
def MainStatement : Prop :=
  ∀ (E : Type) [NormedAddCommGroup E] [NormedSpace ℂ E] [FiniteDimensional ℂ E]
    (X : Type) [TopologicalSpace X] [T2Space X] [CompactSpace X]
    [ChartedSpace E X] [IsManifold 𝓘(ℂ, E) ∞ X],
      letI := HasDerivedCategory.standard (Mod (ringOf (holSheaf E X)))
      (blockFunctor E X).Full ∧ (blockFunctor E X).Faithful ∧
        (blockFunctor E X).essImage = Perf (ringOf (holSheaf E X))
\end{verbatim}
and the theorem, in \texttt{MainH2.lean}:
\begin{verbatim}
theorem main_theorem : Cohesive.MainStatement
-- #print axioms main_theorem : [propext, Classical.choice, Quot.sound]
\end{verbatim}
Reading the statement line by line:
\begin{itemize}
\item \texttt{IsManifold 𝓘(ℂ, E) ∞ X}: $X$ has an atlas with complex-$C^\infty$ (that is,
holomorphic) transition maps. In Mathlib $\infty$ means $C^\infty$; $\top$ would mean
analytic ($C^\omega$). We checked the elaborated statement: the exponent is $\infty$.
\item \texttt{T2Space X}, \texttt{CompactSpace X}: compact Hausdorff; these are implicit in the
sources.
\item \texttt{blockFunctor E X : HoP (dolbeault E X) ⥤ DerivedCategory (Mod (ringOf (holSheaf E X)))}:
the functor $J$ (\S\ref{sec:J}).
\item \texttt{Perf R}: the objects of $D(R)$ isomorphic to $Q(K)$ for a perfect complex $K$
(\S\ref{sec:perf}).
\item \texttt{HasDerivedCategory.standard}: Mathlib's derived category is defined relative to an
instance (a choice of universe and of localization); the statement fixes the standard one,
once, for both \texttt{blockFunctor} and \texttt{Perf}.
\end{itemize}
The conditional form \texttt{MainStatementConditional} (\texttt{StatementsConditional.lean}, frozen
at gate~5) is the same statement with the premises
\texttt{DolbeaultGrothendieck E X → LocalGaugeLemma E X →} after the binders; it is proved in
\texttt{Main.lean} as \texttt{main\_theorem\_conditional}, and \texttt{main\_theorem} applies it to the
proofs \texttt{dolbeaultGrothendieck E X} and \texttt{localGaugeLemma E X}. (Until 6 November 2026
the conditional statement was called \texttt{MainStatement} and the present one
\texttt{MainStatementUnconditional}; the contents did not change.)

\paragraph{Why ``full, faithful, essential image'' and not ``an equivalence onto ${D_{\mathrm{perf}}}$''.}
To write a functor ${\mathrm{Ho}}(\mathcal P_A)\to{D_{\mathrm{perf}}}(X)$ one must already know that $\alpha(M)$ is
perfect, which uses H2 (step C2 below). The functor $J$ is therefore defined into the whole
derived category, without using H1 or H2. The check file \texttt{ChecksMain.lean} proves (K6)
that, for any functor $F$ and any isomorphism-closed property $P$,
\[ F \text{ full} \wedge F\text{ faithful}\wedge \operatorname{essImage}F=P \iff
\exists\, h\colon \forall X,\, P(F X),\ \ (P.\mathrm{lift}\,F\,h) \text{ is an equivalence,} \]
and applies it to \texttt{blockFunctor} and \texttt{Perf} (\texttt{Checks.blockFunctor\_iff}): the
statement is equivalent to ``$J$ corestricts to an equivalence ${\mathrm{Ho}}(\mathcal P_A)\simeq{D_{\mathrm{perf}}}(X)$''.

\paragraph{Sanity checks on the statement.} In complex dimension $0$
($E=\mathrm{Fin}\,0\to\mathbb C$) H1 and H2 are proved directly (K1), independently of their
general proofs, and the point is an example of the theorem (\texttt{blockPerf\_dim\_zero}).
Non-vacuity of every definition (cohesive modules exist, the Dolbeault data satisfy all fields of
the abstract structures, ${\bar\partial}\bar z=d\bar z$ on $\mathbb C$, ${\bar\partial}{\bar\partial}=0$ and Leibniz on
explicit forms on $\mathbb C^2$) is checked in \texttt{Checks*.lean}. Further checks of the
definitions behind the statement: ${\mathrm{Hom}}_{{\mathrm{Ho}}(\mathcal P_A)}({\mathcal{A}},{\mathcal{A}}[n])\cong H^{0,n}_{{\bar\partial}}(X)$
(K4, \S\ref{chap:Z2}), and agreement of our ${\bar\partial}$ with an independent definition by type
projections (K7).

\subsection{The twisted theorem}\label{sec:tw-statement}
For a ${\bar\partial}$-closed $h\in{\mathcal{A}}^{0,2}(X)$, cohesive modules over the \emph{curved} algebra
$({\mathcal{A}}_X(X),{\bar\partial},h)$ have ${\mathbb{X}}\circ{\mathbb{X}}=h\cdot(-)$, so $\alpha(M)$ is no longer a complex. On the
members $U_\alpha$ of a finite cover with $h|_{U_\alpha}={\bar\partial} h^0_\alpha$ the operator
${\mathbb{X}}-h^0_\alpha$ is flat, and the local complexes glue by units $u_{\alpha\beta}$ up to a Čech
2-cocycle ${\check h}$ of units of ${\mathcal{O}}_X$. \cite[Thm~7.4.3]{A} (stated without proof in \cite[Thm~4.2.1]{B})
says that this identifies ${\mathrm{Ho}}(\mathcal P_{(A,{\bar\partial},h)})$ with twisted perfect complexes.
\begin{verbatim}
def MainStatementTwisted : Prop :=
  ∀ (E : Type) [NormedAddCommGroup E] [NormedSpace ℂ E] [FiniteDimensional ℂ E]
    (X : Type) [TopologicalSpace X] [T2Space X] [CompactSpace X]
    [ChartedSpace E X] [IsManifold 𝓘(ℂ, E) ∞ X]
    (h : DolbeaultClosed2 E X) (T : TwistData E X h),
      letI := HasDerivedCategory.standard (TwMod (dolbeaultPreData E X).toSheafData.O T.U T.ch)
      (twistedBlockFunctor h T).Full ∧ (twistedBlockFunctor h T).Faithful ∧
        (twistedBlockFunctor h T).essImage = TwPerf (dolbeaultPreData E X).toSheafData.O T.U T.ch
\end{verbatim}
proved in \texttt{MainTwisted.lean} as \texttt{main\_theorem\_twisted}, with the same three axioms.
Here \texttt{T} is a \emph{twist datum} for $h$ (\S\ref{sec:tw-defs}); one exists for every $h$
(\texttt{exists\_twistData}), so the statement is not vacuous. The check K8
(\texttt{K8\_twisted\_implies\_unconditional}) proves \texttt{MainStatementTwisted → MainStatement}:
at $h=0$ the curved algebra is the flat one by \texttt{rfl}.

%======================================================================
\section{The definitions behind the statements}
The script \texttt{scripts/statement\_defs.py} computes, from the compiled environment, every
project declaration that the statements depend on (types and bodies of definitions; proofs are
not followed since they cannot change the meaning). For \texttt{MainStatement} and
\texttt{MainStatementConditional} the result is 202 declarations in 25 files, frozen at gates 1--5;
with \texttt{MainStatementTwisted} it is 276 declarations in 44 files, all frozen at human review gates,
where they were checked in excerpts (the twisted ones at gate~12). Instances on reached types that are not term dependencies (for
example the triangulated structure on \texttt{HoP}) are listed for information and are not frozen.
This section describes the definitions.

\subsection{Graded-commutative curved dgas (A1)}
A dga is encoded as a $k$-algebra $A$ with a Mathlib \texttt{GradedAlgebra} structure
$\mathcal A\colon\mathbb Z\to\mathrm{Submodule}\,k\,A$ (internal grading) and a class
\texttt{GCDga 𝒜} bundling: ${\mathcal{A}}^q=0$ for $q<0$ and $q>N$, graded commutativity
$ab=(-1)^{pq}ba$, a degree-one derivation $d$ with graded Leibniz rule and $d^2=0$, and a
curvature $c\in{\mathcal{A}}^2$ with $dc=0$ ($d^2=[c,-]$ is automatic, since $c$ is central by graded
commutativity). For Block's theorem $c=0$;
the twisted theorem uses $c=h$. ${\mathcal{A}}^0$ is a commutative ring, and $A$ is an ${\mathcal{A}}^0$-algebra.

\subsection{Cohesive modules and their homotopy category (A3--A5)}
We use \emph{left} modules (Block uses right modules; the translation is the sign convention
$c\leftrightarrow -c$, and a sanity item A13 proves the two definitions dg-equivalent).
\begin{verbatim}
structure QuasiCohesive where
  Q : ℤ → Type                         -- generators Q^p, each an 𝒜⁰-module
  X : Tot 𝒜 Q →ₗ[k] Tot 𝒜 Q          -- 𝕏 on the total module A ⊗_{A⁰} ⨁ Q^p
  X_mem   : ...                        -- 𝕏 has total degree 1
  leibniz : ∀ {p} {a}, a ∈ 𝒜 p → ∀ x,
    X (a • x) = GCDga.d a • x + p.negOnePow • a • X x
  flat    : ∀ x, X (X x) = GCDga.c • x

structure CohesiveModule extends QuasiCohesive 𝒜 where
  bounded    : ∃ a b, ∀ p, (p < a ∨ b < p) → Subsingleton (Q p)
  finite     : ∀ p, Module.Finite (𝒜 0) (Q p)
  projective : ∀ p, Module.Projective (𝒜 0) (Q p)
\end{verbatim}
A cohesive module is thus a bounded graded finitely generated projective ${\mathcal{A}}^0$-module $Q$ with a
$\mathbb Z$-graded superconnection ${\mathbb{X}}$ on $A\otimes_{{\mathcal{A}}^0}Q$ with ${\mathbb{X}}^2=c$ \cite[Def~2.3.2]{B}.

\emph{Hom complexes} (A4): ${\mathrm{Hom}}^n(M_1,M_2)$ consists of $k$-linear maps $\varphi$ of degree $n$
with $\varphi(a\cdot m)=(-1)^{n|a|}a\cdot\varphi(m)$, and
$\delta\varphi={\mathbb{X}}_2\varphi-(-1)^n\varphi{\mathbb{X}}_1$. The sign convention agrees with Mathlib's
\texttt{CochainComplex.HomComplex.δ} (proved). \emph{The homotopy category} \texttt{HoP 𝒜} (A5)
has cohesive modules as objects and $H^0({\mathrm{Hom}}^\bullet(M_1,M_2))$ as morphisms (closed degree-0
maps modulo coboundaries).

\emph{Base change} (A9): for a dga map $f\colon A\to B$, $f^*M$ has generators
$B^0\otimes_{A^0}Q^p$ and ${\mathbb{X}}_B(b\otimes m)=db\otimes m+(-1)^{|b|}b\cdot{\mathbb{X}} m$; it is used to
restrict cohesive modules to open sets. \emph{Modules from complexes} (A10): for $\rho\colon R\to
A^0$ with $d\rho=0$ and a bounded complex $W$ of finitely generated projective $R$-modules,
$A\otimes_R W$ with ${\mathbb{X}}=d\otimes1+1\otimes d_W$ (Koszul sign) is cohesive
(\texttt{CohesiveModule.fromComplex}).

\subsection{\texorpdfstring{Sheaves, perfect complexes, ${D_{\mathrm{perf}}}$ (S1, S9)}{Sheaves, perfect complexes, D\_perf (S1, S9)}}\label{sec:perf}
Sheaves of modules are Mathlib's \texttt{SheafOfModules} over a sheaf of rings
$R$ on \texttt{Opens X}; \texttt{Mod R} abbreviates the category. \texttt{ringOf} turns a sheaf of
commutative rings into a sheaf of rings. Restriction to an open $U$ uses Mathlib's over-site
\texttt{Over U}.
\begin{verbatim}
def IsStrictlyPerfect (P : CochainComplex (Mod R) ℤ) : Prop :=
  (∃ a b : ℤ, P.IsStrictlyGE a ∧ P.IsStrictlyLE b) ∧ ∀ n, IsFinFreeSummand (P.X n)

def IsPerfect (K : CochainComplex (Mod R) ℤ) : Prop :=
  ∀ x : X, ∃ U : Opens X, x ∈ U ∧
    ∃ (L P : CochainComplex (Mod (R.over U)) ℤ) (f : L ⟶ K|U) (g : L ⟶ P),
      QuasiIso f ∧ QuasiIso g ∧ IsStrictlyPerfect P

def Perf (R : RingSheaf X) [HasDerivedCategory.{w} (Mod R)] :
    ObjectProperty (DerivedCategory (Mod R)) :=
  fun D ↦ ∃ K, IsPerfect K ∧ Nonempty (D ≅ DerivedCategory.Q.obj K)
\end{verbatim}
A perfect complex is locally connected by a roof of quasi-isomorphisms to a bounded complex of
direct summands of finite free modules. This is equivalent to the usual
``$Q(K|_U)\cong Q(P)$ in $D(R|_U)$'' but avoids a derived category for every open set.
A sheaf of rings $S$ is \emph{fine} (\texttt{IsFine}) if every finite open cover admits a partition
of unity $\sum\rho_i=1$ by global sections with $\rho_i$ supported in $U_i$ (in a stalk-free
form). On a compact space finite covers suffice.

\subsection{The Dolbeault algebra and holomorphic functions (F1--F12)}
\begin{itemize}
\item \emph{Fibres} (F2): $\Lambda^{0,q}(E;F)$ is the space of continuous $\mathbb C$-alternating
$q$-linear maps on the conjugate space $\bar E$ (\texttt{Conj E}, a type synonym with the conjugate
action). A $\mathbb C$-linear map on $\bar E$ is an antilinear map on $E$, so these are exactly
the $(0,q)$-covectors.
\item \emph{${\bar\partial}$ on open subsets of $E$} (F3): $D''\omega(x)=\tfrac12(D\omega(x)+i\,D\omega(x)\circ J)$,
the antilinear part of the real derivative, and ${\bar\partial}\omega(x)$ is its alternation (Mathlib's
\texttt{alternatizeUncurryFin}). ${\bar\partial}{\bar\partial}=0$ (F4) by copying Mathlib's proof of $d\circ d=0$
with $D''$ for $D$.
\item \emph{Product} (F5): the total fibre $\Lambda^{0,\bullet}=\prod_{q\le n}\Lambda^{0,q}$ gets
its algebra structure transported from Mathlib's exterior algebra of $\mathrm{Dual}_{\mathbb C}\bar
E$ (associativity and graded commutativity for free); Leibniz for ${\bar\partial}$ (F6).
\item \emph{Holomorphic pullback} (F7) commutes with ${\bar\partial}$; this makes the chart-wise definition
independent of charts.
\item \emph{Sheaves on $X$} (F9--F10): a section of ${\mathcal{A}}_X$ over $U$ is a function into the total
fibre, read in charts, whose chart expressions are smooth (a Mathlib \texttt{LocalPredicate}).
${\mathcal{A}}_X(U)$ is a bounded graded-commutative dga with ${\bar\partial}$ (\texttt{dolbeaultGrade E U}); the
Dolbeault algebra is \texttt{dolbeault E X := dolbeaultGrade E ⊤}. The degree-$q$ component sheaves
are subtypes of the total section algebra, so that ${\mathcal{A}}^0$ is literally the dga's degree-0 part.
\item \emph{Holomorphic functions} (F12): \texttt{holSheaf E X := smoothSheafCommRing 𝓘(ℂ, E) 𝓘(ℂ, ℂ) X ℂ},
complex-$C^\infty$ functions, which are the holomorphic ones; $\rho\colon{\mathcal{O}}_X\to{\mathcal{A}}^0$ is the
inclusion.
\end{itemize}

\subsection{The hypotheses H1 and H2}\label{sec:hyp}
They are \texttt{Prop}-valued definitions in \texttt{Hypotheses.lean} (frozen at gate~3). They occur
in \texttt{MainStatementConditional} and, as fields, in the abstract resolving data below; they do
not occur in \texttt{MainStatement}.
\begin{verbatim}
def DolbeaultGrothendieck : Prop :=
  ∀ (U : Opens X) (q : ℕ) (x : X), x ∈ U →
    ∀ α ∈ dolbeaultGrade E U ((q : ℤ) + 1), GCDga.d α = 0 →
      ∃ (U' : Opens X) (h : U' ≤ U), x ∈ U' ∧
        ∃ η ∈ dolbeaultGrade E U' (q : ℤ), GCDga.d η = (dolbeaultRes E h) α

def LocalGaugeLemma : Prop :=
  ∀ (M : CohesiveModule (dolbeault E X)) (x : X),
    ∃ (U' : Opens X) (_ : x ∈ U')
      (W : CochainComplex (ModuleCat (O_X(U'))) ℤ) (hW : IsBoundedFGProj _ W),
      Nonempty (HoP.mk (M.baseChange (dolbeaultRes E (le_top : U' ≤ ⊤))) ≅
                HoP.mk (CohesiveModule.fromComplex (ρ := dolbeaultRhoU E U') W ...))
\end{verbatim}
H1 together with F12 says exactly that $0\to{\mathcal{O}}_X\to{\mathcal{A}}^{0,0}\xrightarrow{{\bar\partial}}{\mathcal{A}}^{0,1}\to\cdots$ is
a resolution. H2 says: near each point, the restriction $M|_{U'}$ is isomorphic in
${\mathrm{Ho}}(\mathcal P_{{\mathcal{A}}(U')})$ to ${\mathcal{A}}(U')\otimes_{{\mathcal{O}}(U')}W$. It is stated for sections over
small opens, only for restrictions of global modules, which is all the proof uses (\cite{A}'s
Assumption~6 is a sheaf version on every open).

\subsection{\texorpdfstring{The functor $J$ (C0, C1)}{The functor J (C0, C1)}}\label{sec:J}
The comparison is proved for abstract \emph{resolving data}, then applied to the Dolbeault data.
\texttt{PreResolvingData X} bundles, for a topological space $X$: a presheaf of bounded
graded-commutative dgas $U\mapsto{\mathcal{A}}(U)$ that is a sheaf (gluing, locality of degree), with flat
curvature; a sheaf of commutative rings ${\mathcal{O}}$ and ring maps $\rho_U\colon{\mathcal{O}}(U)\to{\mathcal{A}}^0(U)$,
natural, with $d\circ\rho=0$; ${\mathcal{A}}^0$ fine; each ${\mathcal{A}}^q$ locally free over ${\mathcal{A}}^0$ (local frames).
\texttt{ResolvingData} adds the resolution condition ($\rho$ injective, exactness in degree $0$ and
in degrees $q+1$) and the local lifting field \texttt{LL} (the abstract form of H2).

For pre-resolving data $P$ and a cohesive module $M$ over $A={\mathcal{A}}(X)$, the complex
$\alpha(M)$ (Lean: \texttt{P.sheafOf M}) has
\[ \alpha(M)^n(V) = \bigl({\mathcal{A}}(V)\otimes_{{\mathcal{A}}^0(V)}\textstyle\bigoplus_p {\mathcal{A}}^0(V)\otimes_{{\mathcal{A}}^0(X)}Q^p\bigr)^n
\;\cong\; \textstyle\bigoplus_{p+q=n}{\mathcal{A}}^q(V)\otimes_{{\mathcal{A}}^0(X)}Q^p, \]
the degree-$n$ part of the total module of the base change $M|_V$, with the differential ${\mathbb{X}}$ of
$M|_V$; ${\mathcal{O}}$ acts through $\rho_V$; restriction maps are the base-change maps. (The first form
reuses the base change A9 and its functoriality; it is canonically isomorphic to Block's
formula.) It is a sheaf (C1.a: $Q^p$ is finitely generated projective and ${\mathcal{A}}^q$ is a summand of
a finite free ${\mathcal{A}}^0$-module). A closed degree-0 map gives a chain map, a coboundary a
null-homotopic one (C1.b), so
\[ J := ({\mathrm{Ho}}(\mathcal P_A)\to K(\mathrm{Mod}\,{\mathcal{O}})\to D(\mathrm{Mod}\,{\mathcal{O}})),\qquad
\texttt{P.J := P.JK ⋙ DerivedCategory.Qh}. \]
Finally (\texttt{BlockFunctorDef.lean}):
\begin{verbatim}
def dolbeaultPreData : PreResolvingData X   -- the Dolbeault data; every field a theorem
def blockFunctor [h : HasDerivedCategory.{w} (Mod (ringOf (holSheaf E X)))] :
    HoP (dolbeault E X) ⥤ DerivedCategory (Mod (ringOf (holSheaf E X))) :=
  @PreResolvingData.J X _ (dolbeaultPreData E X) h
\end{verbatim}
Every field of \texttt{dolbeaultPreData} is a theorem of layer F: gluing and locality of forms,
$\rho$ natural with ${\bar\partial}\rho=0$, ${\mathcal{A}}^0$ fine (smooth partitions of unity, needs compact
Hausdorff), ${\mathcal{A}}^{0,q}$ locally free of rank $\binom nq$. No analytic hypothesis is used:
\texttt{blockFunctor} is defined without H1, H2.

\subsection{Curved Dolbeault algebras, twist data, twisted sheaves (Z3)}\label{sec:tw-defs}
\begin{itemize}
\item \texttt{DolbeaultClosed2 E X}: global forms $h$ of degree $2$ with ${\bar\partial} h=0$;
\texttt{curvedDolbeault h} is $({\mathcal{A}}_X(X),{\bar\partial},h)$, the Dolbeault algebra with curvature replaced
by $h$ (\texttt{GCDga.withC}, on the same carrier).
\item \texttt{TwistData E X h}: a finite open cover $U=(U_\alpha)$, primitives
$h^0_\alpha\in{\mathcal{A}}^{0,1}(U_\alpha)$ with ${\bar\partial} h^0_\alpha=h|_{U_\alpha}$, units
$u_{\alpha\beta}\in{\mathcal{A}}^0(U_\alpha\cap U_\beta)^\times$ with
${\bar\partial} u_{\alpha\beta}=(h^0_\beta-h^0_\alpha)u_{\alpha\beta}$, $u_{\alpha\alpha}=1$,
$u_{\beta\alpha}=u_{\alpha\beta}^{-1}$, and a normalized Čech 2-cocycle ${\check h}$ of units of
${\mathcal{O}}_X$ with $\rho({\check h}_{\alpha\beta\gamma})=u_{\beta\gamma}u_{\alpha\beta}u_{\alpha\gamma}^{-1}$.
\item \texttt{TwMod O U ȟ}: ${\check h}$-twisted sheaves of modules, i.e.\ sheaves $M_\alpha$ on the
over-sites \texttt{Over U\textsubscript{α}} with gluing maps $\varphi_{\alpha\beta}$ on sections over every
$W\le U_\alpha\cap U_\beta$, $\varphi_{\alpha\alpha}=\mathrm{id}$ and
$\varphi_{\beta\gamma}\circ\varphi_{\alpha\beta}={\check h}_{\alpha\beta\gamma}\varphi_{\alpha\gamma}$. It is an
abelian category (its \texttt{Abelian} instance is frozen, since the derived category depends on it)
with enough injectives.
\item \texttt{TwPerf}: objects of $D(\mathrm{TwMod})$ isomorphic to $Q(K)$ with every component
$K_\alpha$ perfect on $U_\alpha$.
\item \texttt{twistedBlockFunctor h T}: $M$ goes to the twisted complex with components
$\alpha(M)^n|_{U_\alpha}$, differential ${\mathbb{X}}-h^0_\alpha$ and gluing $u_{\alpha\beta}\cdot$.
\end{itemize}

%======================================================================
